% P03.1 global-section-palette descent proof, live integration version.

\begin{proof}[\Cref{thm:p031-section-palette-envelope-descent}]
Fix $u$ and $x$.  The finite nonempty fiber has a point $e_x$ at which
$u$ attains its maximum.  Choose $i=i(x,u)$ in
\eqref{eq:p031-global-section-palette} such that $\sigma_i(x)=e_x$, and set
\[
                            b_{x,u}:=u\circ\sigma_i.
\]
Real piecewise linearity is preserved by analytic pullback by
\cite[Lemma~5.4\textup{(1)}, p.~171]{GJR26}, and semipositivity is
preserved by
\cite[Proposition~7.6\textup{(1)}, p.~189]{GJR26}.  Thus $b_{x,u}$ is
semipositive real PL on $X$.  Since $\pi\circ\sigma_i$ is the identity and
$u\leq f\circ\pi$, one has $b_{x,u}\leq f$.  Its pullback is constant on
$\pi^{-1}(x)$, with value
\[
 b_{x,u}(x)=u(\sigma_i(x))=u(e_x)
          =\max_{e'\in\pi^{-1}(x)}u(e').
\]
This proves \eqref{eq:p031-section-palette-exact-repair}.  The index was
selected after $x$ and $u$; no single section is asserted to maximize on
other fibers.

The same two pullback results show that every base competitor pulls back to
a cover competitor.  Hence
\[
                 s_{\mathrm{PL},X}\circ\pi
                    \leq s_{\mathrm{PL},E}.
\]
Conversely, fix $e\in E$, put $x:=\pi(e)$, and let $u$ be a cover
competitor.  Choose $j$ with $\sigma_j(x)=e$.  The preceding argument,
without choosing a maximizing point, makes $u\circ\sigma_j$ a base
competitor.  Therefore
\[
                   u(e)=(u\circ\sigma_j)(x)
                         \leq s_{\mathrm{PL},X}(x).
\]
Taking the supremum over $u$ proves the reverse inequality.  Moreover, a
base competitor pulls back to a cover competitor, while composition with
any palette section sends a cover competitor to a base competitor.  The
two families are therefore empty simultaneously, so the equality includes
the value $-\infty$.

Every locally Fubini--Study function is semipositive real PL by
\cite[Definition~7.8 and \S7.9]{GJR26}.  Adjoining LFS minorants therefore
does not enlarge either real-PL competitor family, and the P03 union
identity follows from
\eqref{eq:p031-section-palette-envelope-pullback}.
\end{proof}

\begin{proof}[\Cref{cor:p031-sectioned-orbit-defect-vanishing}]
For every $g\in G$,
\[
                  \pi\circ g\circ\sigma
                    =\pi\circ\sigma
                    =\operatorname{id}_X,
\]
so $g\circ\sigma$ is a section.  If $e\in\pi^{-1}(x)$, fiber
transitivity gives $e=g(\sigma(x))$ for some $g\in G$.  The translates of
$\sigma$ therefore form a global section palette, with repetitions allowed
when a stabilizer is nontrivial.  The two envelope identities follow from
\Cref{thm:p031-section-palette-envelope-descent}.  Under the stated
nonemptiness hypothesis, subtracting the PL envelope identity in
\eqref{eq:p031-pl-symmetry-defect} gives $\Sigma_{\pi,f}=0$.
\end{proof}

\begin{proof}[\Cref{cor:p031-fold-envelope-without-maximum-closure}]
The two canonical inclusions $X\to X\amalg X$ form a global section
palette for $q$.  Hence
\Cref{thm:p031-section-palette-envelope-descent} gives exact semipositive
real-PL envelope pullback.

Both $h$ and $0$ are base competitors below $f$.  Their pointwise supremum
is $f$, while the envelope is bounded above by the obstacle.  Hence
$s_{\mathrm{PL},X}=f$, and the pullback identity gives
$s_{\mathrm{PL},X\amalg X}=f\circ q$.  The constant zero function is an
LFS minorant of $f$, so the defect is defined and vanishes.

The two values of $u$ on the fiber over $x$ are $h(x)$ and $0$.  Thus
$q_{\max}u=\max\{h,0\}=f$, which is not real PL by
\Cref{prop:p031-real-pl-fiber-maximum-failure}.  The same example therefore
has a failed real-PL maximum operation and exact envelope descent.
\end{proof}
